\section{万有覆盖与覆盖变换}
\sourcepage{31--32}
流形$X$的所有覆盖空间中，有一个可称为“最大”的覆盖，即万有覆盖。所有其他覆盖空间都可作为它的商得到；万有覆盖在“覆盖变换”群作用下的行为，与$X$的基本群密切相关。本节研究这些概念。
\entry{5.1}{定义}
设$X,Y$为连通拓扑空间，$p:Y\to X$为覆盖映射。若满足下述泛性质，则称$p$为$X$的\textbf{万有覆盖}（universal covering）：对每个覆盖$q:Z\to X$（$Z$连通），以及每组满足$p(y_0)=q(z_0)$的点$y_0\in Y,z_0\in Z$，恰存在一个连续保纤维映射$f:Y\to Z$使$f(y_0)=z_0$。

连通拓扑空间$X$的万有覆盖在同构意义下至多一个。事实上，沿用上述记号，若$q:Z\to X$也是万有覆盖，则存在连续保纤维映射$g:Z\to Y$满足$g(z_0)=y_0$。复合$g\circ f:Y\to Y$、$f\circ g:Z\to Z$均连续且保纤维，并分别固定$y_0,z_0$。由泛性质，每种情形满足这些条件的映射只有一个，故$g\circ f=\mathrm{id}_Y$、$f\circ g=\mathrm{id}_Z$。所以$f$是保纤维同胚。
\entry{5.2}{定理}
设$X,Y$为连通流形，$Y$单连通，$p:Y\to X$为覆盖映射。则$p$是$X$的万有覆盖。
\begin{proof}由定义及定理(4.17)直接得到。\end{proof}
\entry{5.3}{定理}
设$X$是连通流形。则存在连通、单连通流形$\widetilde X$及覆盖映射$p:\widetilde X\to X$。

由定理(5.2)，$\widetilde X\to X$即为$X$的万有覆盖。
\begin{proof}
固定$x_0\in X$。对$x\in X$，用$\pi(x_0,x)$表示从$x_0$到$x$的曲线同伦类之集。令
\[
\widetilde X:=\{(x,\alpha):x\in X,\ \alpha\in\pi(x_0,x)\},\qquad p(x,\alpha):=x.
\]
现在在$\widetilde X$上定义拓扑，使其成为连通、单连通的Hausdorff流形，并使$p$为覆盖映射。

设$(x,\alpha)\in\widetilde X$，$U\subset X$是$x$的开、连通、单连通邻域。定义$[U,\alpha]$为所有$(y,\beta)$的集合，其中$y\in U$，且$\beta=\operatorname{cl}(u\cdot v)$，$u$为代表$\alpha$的从$x_0$到$x$的曲线，$v$为完全位于$U$内的从$x$到$y$的曲线。由于$U$单连通，$\beta$与$v$的选择无关。令$\mathfrak B$为所有这些集合所成的族。

断言(a)：$\mathfrak B$是某个拓扑的基。
(i) 每个点显然属于至少一个$[U,\alpha]$。
(ii) 若$(z,\gamma)\in[U,\alpha]\cap[V,\beta]$，则可取$z$的开、连通、单连通邻域$W\subset U\cap V$。容易验证$(z,\gamma)\in[W,\gamma]\subset[U,\alpha]\cap[V,\beta]$。由(i)、(ii)即得断言。

断言(b)：$p$局部同胚，特别地连续。这是因为对每个$[U,\alpha]\in\mathfrak B$，映射$p|_{[U,\alpha]}:[U,\alpha]\to U$为同胚。

断言(c)：$\widetilde X$为Hausdorff空间。
\sourcepage{33}
只须证明同一纤维中不同的两点$(x,\alpha),(x,\beta)$有不交邻域。取$x$的开、连通、单连通邻域$U$。则$[U,\alpha]\cap[U,\beta]=\varnothing$。否则，设$(y,\gamma)$属于交集，选$U$内从$x$到$y$的曲线$w$，并令$\alpha=\operatorname{cl}(u)$、$\beta=\operatorname{cl}(v)$，则按定义$\gamma=\operatorname{cl}(u\cdot w)=\operatorname{cl}(v\cdot w)$。因而$\operatorname{cl}(u)=\operatorname{cl}(v)$，与$\alpha\ne\beta$矛盾。

断言(d)：$\widetilde X$连通，且$p$具有曲线提升性质，从而由(4.19)为覆盖映射。
设$u:[0,1]\to X$以$x_0$为起点。对$s\in[0,1]$定义$u_s(t):=u(st)$，它经过$u$上参数在$[0,s]$中的点。再设$v$为以$x_0$为起终点的闭曲线，则
\[
\widetilde u:[0,1]\to\widetilde X,\quad t\mapsto(u(t),\operatorname{cl}(v\cdot u_t))
\]
连续，是$u$的提升，且$\widetilde u(0)=(x_0,\operatorname{cl}(v))$。这直接由拓扑定义得到。特别地，取$v$为常曲线，便可从$(x_0,e)$到达$\widetilde X$的任意点，故$\widetilde X$连通。最后，若$w$的起点为任意$x_1$，$\alpha\in\pi(x_0,x_1)$，并取从$x_0$到$x_1$且代表$\alpha$的曲线$v$，则$u:=v\cdot w$以$(x_0,e)$为起点的提升给出$w$以$(x_1,\alpha)$为起点的提升；其中$e$为$x_0$处常曲线的同伦类。

断言(e)：$\widetilde X$单连通。
设$w:[0,1]\to\widetilde X$为起终点$(x_0,e)$的闭曲线。则$u:=p\circ w$为$X$中起终点$x_0$的闭曲线。取(d)中$v$为$x_0$处常曲线，得$u$的提升$\widetilde u$。由唯一性，$\widetilde u=w$。故$\widetilde u(1)=(x_0,\operatorname{cl}(u))=(x_0,e)$，所以$u$零伦。由(4.10)，$w$也零伦，因而$\widetilde X$单连通。定理证毕。
\end{proof}
\textbf{注记。}特别地，可以构造任意黎曼曲面的万有覆盖；由(4.6)，它也自然成为黎曼曲面。
\entry{5.4}{定义}
设$p:Y\to X$为拓扑空间之间的覆盖映射。其\textbf{覆盖变换}（covering transformation, deck transformation）是指保纤维同胚$f:Y\to Y$。所有覆盖变换在映射复合下构成群，记为$\operatorname{Deck}(Y/X)$。可能混淆时，写$\operatorname{Deck}(Y\xrightarrow{p}X)$。
\entry{5.5}{定义}
设$X,Y$为连通Hausdorff空间，$p:Y\to X$为覆盖映射。若对每组满足$p(y_0)=p(y_1)$的点$y_0,y_1\in Y$，均存在覆盖变换$f:Y\to Y$使$f(y_0)=y_1$，则称覆盖为\textbf{Galois覆盖}（亦常称正规覆盖或正则覆盖）。
\sourcepage{34}
\textbf{注记。}由(4.8)，满足$f(y_0)=y_1$的覆盖变换至多一个，因为$f$是$p:Y\to X$的提升。

\textbf{例。}覆盖$p:\C^*\to\C^*$，$z\mapsto z^k$为Galois覆盖：若$p(z_1)=p(z_2)$，则$z_2=\omega z_1$，其中$\omega$为$k$次单位根，而$z\mapsto\omega z$是覆盖变换。Galois覆盖与域的Galois扩张存在联系，参见(8.12)。
\entry{5.6}{定理}
设$X$为连通流形，$p:\widetilde X\to X$为万有覆盖。则$p$为Galois覆盖，且$\operatorname{Deck}(\widetilde X/X)$与基本群$\pi_1(X)$同构。
\begin{proof}
(a) 若$p(y_0)=p(y_1)$，由万有覆盖的定义，存在连续保纤维映射$f:\widetilde X\to\widetilde X$使$f(y_0)=y_1$。须证明$f$为同胚。同样存在连续保纤维映射$g$使$g(y_1)=y_0$。则$f\circ g$、$g\circ f$连续且保纤维，分别固定$y_1,y_0$。再次由泛性质，两者均为恒等映射。因此$f$是同胚，即覆盖变换，故$p$为Galois覆盖。

(b) 固定$p(y_0)=x_0$，定义$\Phi:\operatorname{Deck}(\widetilde X/X)\to\pi_1(X,x_0)$如下：若$\sigma$为覆盖变换，选$\widetilde X$中从$y_0$到$\sigma(y_0)$的曲线$v$。由$\widetilde X$单连通，其同伦类唯一确定。曲线$p\circ v$以$x_0$为起终点，令$\Phi(\sigma):=\operatorname{cl}(p\circ v)$。

(i) $\Phi$为群同态。设$v,w$分别从$y_0$到$\sigma(y_0),\tau(y_0)$。则$\sigma\circ w$从$\sigma(y_0)$到$\sigma\tau(y_0)$，且$p\circ(\sigma\circ w)=p\circ w$。乘积$v\cdot(\sigma\circ w)$从$y_0$到$\sigma\tau(y_0)$，故
\begin{align*}
\Phi(\sigma\tau)&=\operatorname{cl}(p\circ(v\cdot(\sigma\circ w)))\\
&=\operatorname{cl}(p\circ v)\operatorname{cl}(p\circ(\sigma\circ w))
=\Phi(\sigma)\Phi(\tau).
\end{align*}
\sourcepage{35}
(ii) $\Phi$单射。若$\Phi(\sigma)=e$，即$p\circ v$零伦，则由(4.10)，提升$v$的终点$\sigma(y_0)$等于起点$y_0$。因此$\sigma=\mathrm{id}_{\widetilde X}$。

(iii) $\Phi$满射。给定$\alpha\in\pi_1(X,x_0)$，选代表曲线$u$，令$v$为它以$y_0$为起点的提升，终点记为$y_1$。由(a)，存在$\sigma\in\operatorname{Deck}(\widetilde X/X)$使$\sigma(y_0)=y_1$。按定义$\Phi(\sigma)=\alpha$。
\end{proof}
\entry{5.7}{例}
(a) $\exp:\C\to\C^*$为万有覆盖，因为$\C$单连通。对$n\in\Z$，令$\tau_n$为平移$2\pi in$。则$\exp(\tau_n(z))=\exp(z+2\pi in)=\exp(z)$，故$\tau_n$为覆盖变换。若$\sigma$为任意覆盖变换，则$\exp(\sigma(0))=1$，故某个$n\in\Z$满足$\sigma(0)=2\pi in=\tau_n(0)$，因而$\sigma=\tau_n$。所以
\[
\operatorname{Deck}(\C\xrightarrow{\exp}\C^*)=\{\tau_n:n\in\Z\}\cong\Z,\qquad\pi_1(\C^*)\cong\Z.
\]
(b) 令$H:=\{z\in\C:\operatorname{Re}z<0\}$为左半平面，$D^*:=\{z\in\C:0<|z|<1\}$。则$\exp:H\to D^*$为穿孔单位圆盘的万有覆盖。与(a)同样可证，覆盖变换群由$2\pi i$整数倍的全部平移组成，$\pi_1(D^*)\cong\Z$。

(c) 若$\Gamma=\Z\gamma_1+\Z\gamma_2$为格，则典范商映射$\C\to\C/\Gamma$为环面的万有覆盖。对$\gamma\in\Gamma$，令$\tau_\gamma$为平移$\gamma$。与(a)类似，覆盖变换群恰为$\{\tau_\gamma:\gamma\in\Gamma\}$，故$\pi_1(\C/\Gamma)\cong\Gamma\cong\Z\times\Z$。

\textbf{推论。}不存在关于$\Gamma$双周期且模$\Gamma$恰有一个一阶极点的亚纯函数。
\begin{proof}
这样的函数将定义全纯映射$f:\C/\Gamma\to\PP^1$，它只取值$\infty$一次。由(4.24)、(2.5)，$f$双全纯，特别地$\pi_1(\C/\Gamma)\cong\pi_1(\PP^1)=0$，矛盾。
\end{proof}
\sourcepage{36}
\textbf{注记。}后面(18.3)将给出具有指定主部的双周期亚纯函数存在的充要条件。但应注意，上述结论仅凭拓扑理由即可得到。
\entry{5.8}{定义}
设$p:Y\to X$为覆盖映射，$G$为$\operatorname{Deck}(Y/X)$的子群。若存在$\sigma\in G$使$\sigma(y)=y'$，则称$y,y'\in Y$模$G$等价。显然这确为$Y$上的等价关系。
\entry{5.9}{定理}
设$X,Y$为连通流形，$q:Y\to X$为覆盖映射，$p:\widetilde X\to X$为万有覆盖。设$f:\widetilde X\to Y$为由万有覆盖定义保证存在的连续保纤维映射。则$f$为覆盖映射，且存在子群$G\subset\operatorname{Deck}(\widetilde X/X)$，使$f(x)=f(x')$当且仅当$x,x'$模$G$等价。此外$G\cong\pi_1(Y)$。
\begin{proof}
先证$f$局部同胚。设$x\in\widetilde X$，$s:=p(x)$，$y:=f(x)$。取$x,s$的开邻域$W_1,U_1$使$p|_{W_1}:W_1\to U_1$为同胚。由$q$为覆盖，可取$s$的连通开邻域$U\subset U_1$，使$q^{-1}(U)=\bigcup_{i\in I}V_i$为不交开集之并，各$q|_{V_i}:V_i\to U$为同胚。令$V$为包含$y$的那一个$V_i$，$W:=p^{-1}(U)\cap W_1$。由$y\in f(W)\subset q^{-1}(U)$及$f(W)$连通，得$f(W)=V$。由于$p|_W$及$q|_V$均为到$U$的同胚，$f|_W:W\to V$也为同胚。

为证$f$为覆盖，考虑$Y$中以$y_0$为起点的曲线$v$，以及$f(x_0)=y_0$的点$x_0$。曲线$q\circ v$可关于$p$提升为以$x_0$为起点的曲线$u$。则$f\circ u$与$v$均为$q\circ v$的提升，且起点同为$y_0$，故相等。因此$u$正是$v$关于$f$的所需提升。由(4.19)，$f$为覆盖。

令$G:=\operatorname{Deck}(\widetilde X/Y)$，它是$\operatorname{Deck}(\widetilde X/X)$的子群。由于$\widetilde X$单连通，$f$是$Y$的万有覆盖，故为Galois覆盖。因此$G\cong\pi_1(Y)$，且$f(x)=f(x')$当且仅当存在$\sigma\in G$使$\sigma(x)=x'$。
\end{proof}
现在用(5.9)确定穿孔单位圆盘$D^*$的所有覆盖空间。
\entry{5.10}{定理}
设$X$是黎曼曲面，$f:X\to D^*$为无分歧全纯覆盖映射。则：
\sourcepage{37}
(i) 若覆盖叶数无限，则存在$X$到左半平面$H$的双全纯映射$\varphi$，满足
\begin{equation}\exp\circ\varphi=f.\tag{1}\end{equation}
(ii) 若覆盖为$k$叶（$k<\infty$），则存在双全纯映射$\varphi:X\to D^*$满足
\begin{equation}p_k\circ\varphi=f,\qquad p_k(z)=z^k.\tag{2}\end{equation}
这些等式分别表达原书交换三角图(1)、(2)中的关系。因此，每个$D^*$的覆盖，都同构于由对数或$k$次根给出的覆盖。
\begin{proof}
由于$\exp:H\to D^*$为万有覆盖，存在全纯映射$\psi:H\to X$满足$\exp=f\circ\psi$。令$G\subset\operatorname{Deck}(H/D^*)$为对应子群。

(i) 若$G$只含单位元，则$\psi:H\to X$双全纯，所求$\varphi$为其逆。

(ii) $\operatorname{Deck}(H/D^*)=\{\tau_n:n\in\Z\}$，其中$\tau_n(z)=z+2\pi in$。因此每个非平凡子群$G$都形如$\{\tau_{nk}:n\in\Z\}$，其中$k\ge1$。令$g:H\to D^*$，$g(z):=\exp(z/k)$。则$g(z)=g(z')$当且仅当$z,z'$模$G$等价。因此存在双射$\varphi:X\to D^*$使$\varphi\circ\psi=g$，即原书相应三角图交换。由于$\psi,g$局部双全纯，$\varphi$双全纯。容易验证(2)，定理得证。
\end{proof}
\sourcepage{38}
\entry{5.11}{定理}
设$X$是黎曼曲面，$D$为单位圆盘，$f:X\to D$为固有非常值全纯映射，且在$D^*=D\setminus\{0\}$之上无分歧。则存在自然数$k\ge1$及双全纯映射$\varphi:X\to D$，使
\begin{equation}p_k\circ\varphi=f,\qquad p_k(z):=z^k.\tag{*}\end{equation}
\begin{proof}
令$X^*:=f^{-1}(D^*)$，则$f|_{X^*}:X^*\to D^*$是无分歧固有全纯覆盖。由上一定理，存在双全纯映射$\varphi:X^*\to D^*$使$p_k\circ\varphi=f|_{X^*}$。断言$f^{-1}(0)$只有一点。反设它含$n\ge2$个点$b_1,\ldots,b_n$。则存在两两不交的$b_i$的开邻域$V_i$及圆盘$D(r):=\{z:|z|<r\}$（$0<r\le1$），满足
\begin{equation}f^{-1}(D(r))\subset V_1\cup\cdots\cup V_n.\tag{**}\end{equation}
令$D^*(r)=D(r)\setminus\{0\}$。由于$f^{-1}(D^*(r))$同胚于$p_k^{-1}(D^*(r))=D^*(\sqrt[k]r)$，它连通。又因每个$b_i$都是$f^{-1}(D^*(r))$的聚点，$f^{-1}(D(r))$也连通，与(**)矛盾。因此$f^{-1}(0)=\{b\}$。定义$\varphi(b):=0$，便可把$\varphi$延拓为双全纯映射$X\to D$，且满足(*)。
\end{proof}
\subsection*{习题（第5节）}
\entry{5.1}{}
令$X=\C\setminus\{\pm1\}$，$Y=\C\setminus\{\pi/2+k\pi:k\in\Z\}$（参见习题4.1）。证明$\operatorname{Deck}(Y\xrightarrow{\sin}X)$由以下变换组成：
(i) $f_k(z)=z+2k\pi$；(ii) $g_k(z)=-z+(2k+1)\pi$，其中$k\in\Z$。计算乘积$f_k\circ f_l$、$f_k\circ g_l$、$g_l\circ f_k$、$g_k\circ g_l$。
\sourcepage{39}
\entry{5.2}{}
设$X$是连通流形，$p:\widetilde X\to X$为万有覆盖。取子群$G\subset\operatorname{Deck}(\widetilde X/X)$，令$Y:=\widetilde X/G$为按(5.8)中等价关系所得商空间，$q:Y\to X$为$p$诱导的映射。证明$q$为覆盖映射，且为Galois覆盖当且仅当$G$为$\operatorname{Deck}(\widetilde X/X)$的正规子群；在后一情形，$\operatorname{Deck}(Y/X)\cong\operatorname{Deck}(\widetilde X/X)/G$。
\entry{5.3}{}
求覆盖$\tan:\C\to\PP^1\setminus\{i,-i\}$的覆盖变换（参见习题4.4）。
\entry{5.4}{}
设$\Gamma,\Gamma'\subset\C$为格，$f:\C/\Gamma\to\C/\Gamma'$为非常值全纯映射，且$f(0)=0$。证明存在唯一$\alpha\in\C^*$使$\alpha\Gamma\subset\Gamma'$，且有交换图
\[
\begin{array}{ccc}\C&\xrightarrow{F}&\C\\ {\scriptstyle\pi}\downarrow&&\downarrow{\scriptstyle\pi'}\\ \C/\Gamma&\xrightarrow{f}&\C/\Gamma'\end{array}
\]
其中$F(z)=\alpha z$，$\pi,\pi'$为典范投影。证明$f$为无分歧覆盖，且$\operatorname{Deck}(\C/\Gamma\xrightarrow{f}\C/\Gamma')\cong\Gamma'/\alpha\Gamma$。
\entry{5.5}{}
令$X:=\C\setminus\{2,-2\}$，$Y:=\C\setminus\{\pm1,\pm2\}$，$p:Y\to X$由$p(z):=z^3-3z$给出。证明$p$为无分歧3叶全纯覆盖。计算$\operatorname{Deck}(Y/X)$并证明此覆盖非Galois。［提示：利用每个双全纯映射$f:Y\to Y$均可延拓为$\PP^1$的自同构这一事实。］
\entry{5.6}{}
令$X:=\C\setminus\{0,1\}$，$Y:=\C\setminus\{0,\pm i,\pm i\sqrt2\}$，$p:Y\to X$由$p(z):=(z^2+1)^2$给出。证明$p$为非Galois的无分歧4叶覆盖，并且$\operatorname{Deck}(Y/X)=\{\mathrm{id},\varphi\}$，其中$\varphi(z):=-z$。
\entry{5.7}{}
设$X,Y$为连通Hausdorff空间。证明每个2叶覆盖$p:Y\to X$都是Galois覆盖。
